\section*{Hoofdstuk \ref{ch:b1:spherical-astronomy} --- Sferische astronomie}

\begin{solution}{exo:b1:spherical-astronomy:1}
De \omterm{def:b1:spherical-astronomy:horizontal}{hoogte} van de pool is gelijk aan de \omterm{def:b1:spherical-astronomy:geographic}{breedtegraad}: \ang{52}. Een ster gaat
nooit onder als $\delta > \ang{90} - \ang{52} = \ang{38}$, en komt nooit op als
$\delta < \ang{-38}$.
\end{solution}

\begin{solution}{exo:b1:spherical-astronomy:2}
$\qty{6.75}{h}\times\ang{15} = \ang{101.25}$. $H = \mathrm{LST} - \alpha =
\qty{10}{h} - \qty{6}{h}\,\qty{45}{min} = \qty{3}{h}\,\qty{15}{min}$, dat is
\ang{48.75}. De \omterm{def:b1:spherical-astronomy:equatorial}{uurhoek} is positief: de ster staat west van de meridiaan, die
zij drie en een kwart sterrenuur geleden passeerde.
\end{solution}

\begin{solution}{exo:b1:spherical-astronomy:3}
$h = \ang{90} - \varphi + \delta_\odot$: \ang{73.44} op de zomerzonnewende,
\ang{50} op de equinoxen, \ang{26.56} op de winterzonnewende.
\end{solution}

\begin{solution}{exo:b1:spherical-astronomy:4}
1900: deelbaar door 100 maar niet door 400, dus geen schrikkeljaar. 2000:
deelbaar door 400, dus een schrikkeljaar. 2100: geen schrikkeljaar. Het
Juliaanse jaar overschrijdt het \omterm{def:b1:spherical-astronomy:tropical-year}{tropisch jaar} met $365.25 - 365.2422 = 0.0078$
dagen, \num{0.78} dagen per eeuw; het Gregoriaanse jaar met $0.0003$ dagen,
\num{0.03} dagen per eeuw.
\end{solution}

\begin{solution}{exo:b1:spherical-astronomy:5}
$\cos\Delta = \sin^2\ang{45} + \cos^2\ang{45}\cos\ang{90} = 0.5$, dus
$\Delta = \ang{60}$ en de afstand is $6371\times\pi/3 = \qty{6672}{km}$. Langs de
parallel, een \omterm{def:b1:spherical-astronomy:spherical-triangle}{kleine cirkel} met straal $R\cos\ang{45}$, is de afstand
$6371\times\cos\ang{45}\times\pi/2 = \qty{7076}{km}$. De \omterm{def:b1:spherical-astronomy:spherical-triangle}{grootcirkel} is ongeveer
\qty{404}{km} korter.
\end{solution}

\begin{solution}{exo:b1:spherical-astronomy:6}
$\sin h = \sin\ang{48}\sin\ang{20} + \cos\ang{48}\cos\ang{20}\cos\ang{30}
= 0.2542 + 0.5446 = 0.7987$, $h = \ang{53.0}$.
$\cos\mathrm{A} = (\sin\ang{20} - \sin\ang{48}\sin\ang{53.0})/(\cos\ang{48}
\cos\ang{53.0}) = (0.3420 - 0.5935)/0.4027 = -0.6246$, $\arccos = \ang{128.7}$;
de ster staat west van de meridiaan ($H > 0$), dus $\mathrm{A} = \ang{360} -
\ang{128.7} = \ang{231.3}$: naar het zuidwesten, \ang{53} boven de horizon.
\end{solution}

\begin{solution}{exo:b1:spherical-astronomy:7}
$\cos H_0 = -\tan\ang{60}\tan\ang{23.44} = -1.7321\times0.4335 = -0.7510$,
$H_0 = \ang{138.67}$, daglengte $2\times138.67/15 = \qty{18.49}{h}$, ongeveer
\qty{18}{h}~\qty{29}{min}, zoals de kromme voor \ang{60} N op de zomerzonnewende
laat zien.
\end{solution}

\begin{solution}{exo:b1:spherical-astronomy:8}
De doorgang schuift elke dag \qty{3.93}{min} vroeger, dus in 30 dagen
$30\times3.93 =
\qty{118}{min}$: de ster kruist om ongeveer 20:02 uur. Om op het
middaguur te kruisen moet zij \qty{10}{h} = \qty{600}{min} vroeger
verschuiven, wat $600/3.93 \approx 153$ dagen kost.
\end{solution}

\begin{solution}{exo:b1:spherical-astronomy:9}
(a) Op een equinox is de middaghoogte \ang{45}: $I = 1361\times\sin\ang{45}
= \qty{962}{W.m^{-2}}$. (b) Op de zomerzonnewende van juni is $h = \ang{68.44}$
en $(a/r)^2 = (149.598/152.100)^2 = 0.9674$:
$I = 1361\times0.9674\times\sin\ang{68.44} = \qty{1224}{W.m^{-2}}$. De hogere
zon compenseert ruimschoots de grotere afstand.
\end{solution}

\begin{solution}{exo:b1:spherical-astronomy:10}
$\sin\delta_\odot = \sin\ang{23.44}\sin\ang{60} = 0.3445$,
$\delta_\odot = \ang{20.15}$. $\tan\alpha_\odot = \cos\ang{23.44}\tan\ang{60}
= 1.5890$, $\alpha_\odot = \ang{57.82}$ (zelfde kwadrant als $\lambda_\odot$).
De \omterm{def:b1:spherical-astronomy:equatorial}{rechte klimming} blijft achter met \ang{2.18}, dat is $2.18\times4 \approx
\qty{8.7}{min}$ tijd. Een zon die gelijkmatig langs de \omterm{def:b1:spherical-astronomy:ecliptic}{ecliptica} beweegt, zou
toch ongelijkmatig bewegen in \omterm{def:b1:spherical-astronomy:equatorial}{rechte klimming}, de grootheid die de \omterm{def:b1:spherical-astronomy:equatorial}{uurhoek} en
dus de zonnewijzer bepaalt: achter de gemiddelde zon tussen een equinox en de
volgende zonnewende, voor haar tussen een zonnewende en de volgende equinox.
Dit is de schuinte-term van de \omterm{def:b1:spherical-astronomy:solar-time}{tijdsvereffening}.
\end{solution}

\begin{solution}{exo:b1:spherical-astronomy:11}
Bij de opkomst geldt $h = 0$, dus \cref{prop:b1:spherical-astronomy:transform}
geeft $\sin\delta = \cos\varphi\cos\mathrm{A}_0$; op de zomerzonnewende geldt
$\delta = \varepsilon$, dus $\cos\mathrm{A}_0 = \sin\varepsilon/\cos\varphi$.
Voor $\mathrm{A}_0 = \ang{45}$: $\cos\varphi = \sin\ang{23.44}/\cos\ang{45} =
0.5626$, $\varphi = \ang{55.8}$ N. Boven $\ang{90} - \varepsilon =
\ang{66.56}$ is $\sin\varepsilon/\cos\varphi > 1$: er is dan geen \omterm{def:b1:spherical-astronomy:horizontal}{azimut} van
opkomst, omdat de zon die dag niet ondergaat (middernachtzon).
\end{solution}

\begin{solution}{exo:b1:spherical-astronomy:12}
(a) Een tweehoek met hoek $\theta$ is het deel $\theta/2\pi$ van de bol:
$\frac{\theta}{2\pi}\times4\pi R^2 = 2\theta R^2$. (b) De drie paren tweehoeken
hebben samen oppervlakte $2\times(2A + 2B + 2C)R^2$; zij bedekken de bol één
keer, behalve de driehoek en haar tegenoverliggende, die elk drie keer worden
bedekt, dus dubbel zoveel als nodig: $4(A + B + C)R^2 = 4\pi R^2 + 4S$, waarin
$S$ de oppervlakte van de driehoek is. Daaruit volgt $S = R^2(A + B + C - \pi)$.
(c) Drie rechte hoeken: $S = R^2(3\pi/2 - \pi) = \pi R^2/2$, een achtste van de
bol (een octant). De driehoek uit de figuur: $89.9 + 64.9 + 67.4 - 180 = \ang{42.2}$, dat is \num{0.736} radialen: $S = 0.736\,R^2$, \qty{5.9}{\percent}
van de bol.
\end{solution}

\begin{solution}{pb:b1:spherical-astronomy:1}
\textbf{1.} Poolafstanden $\ang{90} - (\ang{-24.63}) = \ang{114.63}$ en
$\ang{90} - \ang{51.48} = \ang{38.52}$; lengteverschil \ang{70.40}.

\textbf{2.} $\cos\Delta = \cos\ang{114.63}\cos\ang{38.52} +
\sin\ang{114.63}\sin\ang{38.52}\cos\ang{70.40} = -0.3260 + 0.1899 =
-0.1362$.

\textbf{3.} $\Delta = \ang{97.83}$; afstand $6371\times1.7074 =
\qty{10878}{km}$, ongeveer \qty{10900}{km}.

\textbf{4.}
\[
  \cos\mathrm{Az} = \frac{\cos\ang{38.52} - \cos\ang{114.63}\cos\ang{97.83}}
  {\sin\ang{114.63}\sin\ang{97.83}} = \frac{0.7823 - 0.0568}{0.9006} = 0.8056 ,
\]
$\mathrm{Az} = \ang{36.3}$, naar het noordoosten.

\textbf{5.} $\sin\mathrm{Az} = \sin\ang{70.40}\sin\ang{38.52}/\sin\ang{97.83}
= 0.5922$, wat \ang{36.3} of \ang{143.7} oplevert; de positieve cosinus van
vraag 4 kiest \ang{36.3}.

\textbf{6.} Nee. Een constante koers snijdt elke meridiaan onder dezelfde hoek
en spiraalt naar een pool; de koers langs de \omterm{def:b1:spherical-astronomy:spherical-triangle}{grootcirkel} verandert voortdurend
(zij is alleen aan het begin \ang{36.3}). Het vliegtuig zou een langere weg
afleggen en elders aankomen, tenzij het zijn koers corrigeert.

\textbf{7.} De \omterm{def:b1:spherical-astronomy:horizontal}{hoogte} van de zuidelijke \omterm{def:b1:spherical-astronomy:celestial-sphere}{hemelpool} is gelijk aan de absolute
waarde van de \omterm{def:b1:spherical-astronomy:geographic}{breedtegraad}: \ang{24.63}.

\textbf{8.} $h = \ang{90} - |\ang{-24.63} - (\ang{-16.72})| = \ang{82.09}$;
$\delta > \varphi$, dus Sirius gaat \ang{7.91} ten noorden van het \omterm{def:b1:spherical-astronomy:celestial-sphere}{zenit} door.

\textbf{9.} Een ster gaat nooit onder als $\delta < -(\ang{90} - \ang{24.63}) =
\ang{-65.37}$. Canopus, op \ang{-52.70}, is niet circumpolair: zij gaat onder.

\textbf{10.} $\ang{90} - |\ang{-24.63} + \ang{52.70}| = \ang{61.93}$, ten zuiden
van het \omterm{def:b1:spherical-astronomy:celestial-sphere}{zenit}.

\textbf{11.} $\cos H_0 = -\tan(\ang{-24.63})\tan(\ang{-16.72}) =
-(0.4584)(0.3004) = -0.1377$; $H_0 = \ang{97.92}$. Sirius blijft
$2\times97.92/15 = \qty{13.06}{h}$ \omterm{def:b1:spherical-astronomy:sidereal-time}{sterrentijd} boven de horizon, ongeveer
\qty{13}{h}.

\textbf{12.} $\cos\mathrm{A}_0 = \sin(\ang{-16.72})/\cos\ang{24.63} =
-0.3165$: Sirius komt op bij het \omterm{def:b1:spherical-astronomy:horizontal}{azimut} \ang{108.5}, ten zuiden van het oosten.

\textbf{13.} Wanneer $H = 0$, is $\mathrm{LST} = \alpha = \qty{6}{h}~\qty{45}{min}~\qty{9}{s}$.

\textbf{14.} $\qty{86400}{s}\times365.2422/366.2422 = \qty{86164.09}{s} =
\qty{23}{h}~\qty{56}{min}~\qty{4}{s}$; de passage schuift elke nacht
\qty{235.9}{s}, dus \qty{3}{min}~\qty{56}{s}, vroeger.

\textbf{15.} Nodig is $12 + n\times24/365.2422 = 6.7525 + 24$, dus
$n = 18.7525/0.065710 = 285.4$ dagen. 285 dagen na 20 maart is 30 december:
Sirius kruist de meridiaan om middernacht van de plaatselijke gemiddelde tijd in
de nacht van 30 op 31 december 2026.

\textbf{16.} Paranal ligt \ang{70.40} west van Greenwich: de plaatselijke
gemiddelde tijd loopt $70.40/15 = \qty{4.693}{h}$ achter, \qty{4}{h}~\qty{41.6}{min}.

\textbf{17.} \omterm{def:b1:spherical-astronomy:solar-time}{Ware zonnetijd} = \omterm{def:b1:spherical-astronomy:solar-time}{gemiddelde zonnetijd} + \omterm{def:b1:spherical-astronomy:solar-time}{tijdsvereffening} =
$\qty{0}{h} - \qty{3}{min}$: de zonnewijzer wijst ongeveer 23:57.

\textbf{18.} $7\times\qty{3.93}{min} \approx \qty{27.5}{min}$ eerder.

\textbf{19.} $+\ang{23.44}$ op de zomerzonnewende van juni, $-\ang{23.44}$ op de
winterzonnewende van december.

\textbf{20.} Juni: $\ang{90} - |\ang{-24.63} - \ang{23.44}| = \ang{41.93}$.
December: $\ang{90} - |\ang{-24.63} + \ang{23.44}| = \ang{88.81}$; omdat
$\delta_\odot = \ang{-23.44} > \varphi$, gaat de zon \ang{1.19} ten noorden van
het \omterm{def:b1:spherical-astronomy:celestial-sphere}{zenit} door.

\textbf{21.} $I = 1361\times(1/0.9837)^2\times\sin\ang{88.81} = 1361\times
1.0334\times0.9998 = \qty{1406}{W.m^{-2}}$.

\textbf{22.} In juni, de zuidelijke winter. $\cos H_0 = -\tan(\ang{-24.63})
\tan\ang{23.44} = 0.4584\times0.4335 = 0.1988$: de zon gaat onder bij de
\omterm{def:b1:spherical-astronomy:equatorial}{uurhoek} $H_0 = \ang{78.53}$.

\textbf{23.} $\cos\mathrm{A}_0 = \sin\ang{23.44}/\cos\ang{24.63} = 0.4376$:
$\mathrm{A}_0 = \ang{64.0}$, ten noorden van het oosten.

\textbf{24.} De daglengte is $2\times78.53/15 = \qty{10.47}{h}$; de langste nacht
van het jaar op Paranal duurt $24 - 10.47 = \qty{13.53}{h}$, dat is
\textbf{\qty[text-series-to-math=true]{13}{h}~\qty[text-series-to-math=true]{32}{min}}.
\end{solution}
